%==============================================================================
% Architecture Scaling Ablation — Beamer template deck
%
% Build:   latexmk -pdf main.tex      (or)   pdflatex main.tex
% Figures live in ../figures (see \graphicspath below).
%==============================================================================
\documentclass[aspectratio=169,11pt]{beamer}

%------------------------------------------------------------------------------
% Theme & appearance
%------------------------------------------------------------------------------
\usetheme{default}
\usecolortheme{seahorse}
\setbeamertemplate{navigation symbols}{}      % drop the nav clutter
\setbeamertemplate{caption}[numbered]
\setbeamertemplate{itemize items}[circle]
\setbeamerfont{title}{series=\bfseries}
\setbeamerfont{frametitle}{series=\bfseries}

%------------------------------------------------------------------------------
% Packages
%------------------------------------------------------------------------------
\usepackage[T1]{fontenc}
\usepackage{graphicx}
\usepackage{booktabs}
\usepackage{amsmath}
\usepackage{amssymb}

\graphicspath{{../figures/}{../figures/scaling/}}

% Footer with author / short title / slide number
\setbeamertemplate{footline}{%
  \leavevmode%
  \hbox{%
    \begin{beamercolorbox}[wd=.4\paperwidth,ht=2.25ex,dp=1ex,leftskip=1em]{author in head/foot}%
      \usebeamerfont{author in head/foot}\insertshortauthor%
    \end{beamercolorbox}%
    \begin{beamercolorbox}[wd=.45\paperwidth,ht=2.25ex,dp=1ex,center]{title in head/foot}%
      \usebeamerfont{title in head/foot}\insertshorttitle%
    \end{beamercolorbox}%
    \begin{beamercolorbox}[wd=.15\paperwidth,ht=2.25ex,dp=1ex,rightskip=1em,right]{date in head/foot}%
      \usebeamerfont{date in head/foot}\insertframenumber{}\,/\,\inserttotalframenumber%
    \end{beamercolorbox}%
  }%
  \vskip0pt%
}

%------------------------------------------------------------------------------
% Title metadata
%------------------------------------------------------------------------------
\title[MetaNCA architecture scaling]{MetaNCA architecture scaling}
\subtitle{Capacity, depth, and held-out architecture generalization}
\author[D.~Berenberg]{Dan Berenberg}
\institute{Mythos Scientific}
\date{\today}

%==============================================================================
\begin{document}

%------------------------------------------------------------------------------
% 1. Title slide
%------------------------------------------------------------------------------
\begin{frame}[plain]
  \titlepage
\end{frame}

%------------------------------------------------------------------------------
% 2. Bulleted text template
%------------------------------------------------------------------------------
\begin{frame}{General questions for this study}
  \framesubtitle{Characterizing the scaling laws of MetaNCA}

  \begin{itemize}
    \item Weight transformer should ``scale'' similar to language models
      \begin{itemize}
        \item Text tokens $\rightarrow$ MetaNCA cells $=$ task network parameters
        \item Text documents $\rightarrow$ task network architectures
      \end{itemize}
    \item How does scaling $N_\mathrm{arch}$ or $N_\mathrm{params}$ impact downstream performance?
        \begin{itemize}
            \item Validation loss, accuracy, perplexity, bpb, \ldots
            \item Generalization to unseen architectures: out-of-distribution extrapolation,
                in-distribution interpolation
            \item Impact of scaling on training arch. performance
        \end{itemize}
  \end{itemize}
  %\vfill
  %\begin{block}{Takeaway}
  %  Use a \alert{block} to highlight the one thing the audience should remember.
  %\end{block}
\end{frame}

\begin{frame}{Other questions, not part of this study}

  \begin{itemize}
    \item Understanding/fixing instability during training
    \item Characterizing sample pooling
        \begin{itemize}
            \item Number of update steps in this setting
        \end{itemize}
    \item Update step increment schedule during training
    \item Other scaling laws for the local rule:
        \begin{itemize}
            \item \# of parameters
            \item \# attn. heads 
            \item hidden state dize
            \item architecture / capacity of post-attn. layers
            \item scaling \# of tasks
        \end{itemize}
  \end{itemize}
  %\vfill
  %\begin{block}{Takeaway}
  %  Use a \alert{block} to highlight the one thing the audience should remember.
  %\end{block}
\end{frame}

\begin{frame}{FashionMNIST}

  \begin{itemize}
    \item Small dataset (same size/dimensions as MNIST)
    \item Train MetaNCA with 1, 2, 4, and 8 architectures
    \item Task networks: 
        \begin{itemize}
            \item All MLPs
            \item Fixed depth $\ell = 3$ hidden layers
            \item Nonincreasing \# neurons per layer:
                $N_\mathrm{neurons} \in \{16, 32, 64, 128 \}$
        \end{itemize}
    \item Training:
        \begin{itemize}
            \item Learning rate 0.001
            \item Train for 1.2K metaepochs
            \item Increment \# update steps every 100 metaepochs until \# updates = 10
        \end{itemize}
    \item Validation: 8 validation architectures, tested on val. set
  \end{itemize}
\end{frame}

%%------------------------------------------------------------------------------
%% 3. Left image, right text template
%%------------------------------------------------------------------------------
%\begin{frame}{Left image, right text template}
%  \begin{columns}[T,onlytextwidth]
%    % --- left: figure ---
%    \begin{column}{0.55\textwidth}
%      \centering
%      \includegraphics[width=\linewidth,height=0.75\textheight,keepaspectratio]%
%        {fashion_mnist_capacity_x_perf.pdf}
%      \\[0.3em]
%      {\footnotesize Capacity vs.\ performance.}
%    \end{column}
%    % --- right: text ---
%    \begin{column}{0.42\textwidth}
%      \textbf{Reading the figure}
%      \begin{itemize}
%        \item What the x-axis encodes.
%        \item What the trend shows.
%        \item Why it matters for scaling.
%      \end{itemize}
%      \vfill
%      \emph{Swap the filename above for any figure in \texttt{../figures}.}
%    \end{column}
%  \end{columns}
%\end{frame}

%------------------------------------------------------------------------------
% 4. Right image, left text template
%------------------------------------------------------------------------------
\begin{frame}{FashionMNIST}
  \framesubtitle{Effect of multiple architectures}
  \begin{columns}[T,onlytextwidth]
    % --- left: text ---
    \begin{column}{0.42\textwidth}
      \begin{itemize}
        \item Val. arch. performance scales with the number of training archs
        \item Saturates around Adam performance
      \end{itemize}
    \end{column}
    % --- right: figure ---
    \begin{column}{0.55\textwidth}
      \centering
      \includegraphics[width=\linewidth,height=0.75\textheight,keepaspectratio]%
        {fashion_mnist_fixed_depth.pdf}
      \\[0.3em]
      {\footnotesize}
    \end{column}
  \end{columns}
\end{frame}

\begin{frame}{FashionMNIST}
  \framesubtitle{Distance to training set}
  \begin{columns}[T,onlytextwidth]
    % --- left: text ---
    \begin{column}{0.45\textwidth}
      \textbf{Calculating distance}
      \begin{itemize}
          \item Each model is a depth $\ell = 4$ MLP with weight matrices $\{ W_i \}_{i=1}^\ell$,
              where $W_i \in \mathbb{R}^{d_{i-1} \times d_i}$
            \item $d_0 = 784, d_4 = 10$
        \item Each MLP is encoded by the vector of its weights' codomains, omitting $d_0$ and $d_4$: $\begin{pmatrix} d_1 & d_2 & d_3 \end{pmatrix}$
        \item The distance between two MLPs $A$ and $B$ can be calculated by the cosine distance of their architecture vectors.
      \end{itemize}
      %\vfill
      %\begin{block}{Note}
      %  Mirror of the previous layout — text left, figure right.
      %\end{block}
    \end{column}
    % --- right: figure ---
    \begin{column}{0.5\textwidth}
      \centering
      \includegraphics[width=\linewidth,height=0.75\textheight,keepaspectratio]%
        {fashion_mnist_heldout_arch_perf.pdf}
      \\[0.3em]
      {\footnotesize}
    \end{column}
  \end{columns}
\end{frame}

\begin{frame}{FashionMNIST}
  \framesubtitle{Capacity vs. Performance}
    \begin{itemize}
    \item Correlation between \# parameters and downstream performance: strong in train, weaker in val.
    \end{itemize}
    \centering
    \includegraphics[width=\linewidth,height=0.75\textheight,keepaspectratio]%
    {fashion_mnist_capacity_x_perf.pdf}
\end{frame}

\begin{frame}{Shakespeare}
    \begin{columns}[T, onlytextwidth]
    \begin{column}{0.45\textwidth}
        \begin{itemize}
            \item Architecture scaling for language modeling
            \item 1, 2, 4, and 8 architectures: (3) layer transformers
            \item $d_\mathrm{model} \in \{ n \in \mathbb{N} \mid 32 \leq n \leq 224, \ n \ \mathrm{mod} \ 8 = 0\}$
            \item fixed number of heads (4), feedforward multiplier is ($4\times$), vocab size is 10k 
            \item Train for 330 metaepochs, increment every 30 up to 10 update steps.
            \item learning rate 3e-4
        \end{itemize}
    \end{column}
        \begin{column}{0.5\textwidth}
        \includegraphics[width=\linewidth,height=0.75\textheight,keepaspectratio]%
        {llmw_train_curves_vs_adam.pdf}
        \end{column}
    \end{columns}
    \vfill
    \begin{block}{Takeaway}
        The local rule may be underpowered. Currently investigating.
    \end{block}

\end{frame}

\end{document}
